% tip to kill the wrong type
%!TEX program = xelatex
\documentclass[10pt,journal,compsoc]{IEEEtran}
\IEEEoverridecommandlockouts
% The preceding line is only needed to identify funding in the first footnote. If that is unneeded, please comment it out.
\usepackage{cite}
\usepackage{amsmath,amssymb,amsfonts}
\usepackage{algorithmic}
\usepackage{graphicx}
\usepackage{textcomp}
\usepackage{xcolor}
\usepackage{zh_CN-Adobefonts_external}
\usepackage{algorithm}
\usepackage{xeCJK}
\def\BibTeX{{\rm B\kern-.05em{\sc i\kern-.025em b}\kern-.08em
    T\kern-.1667em\lower.7ex\hbox{E}\kern-.125emX}}
\usepackage{fancyhdr}
%\pagenumbering{roman}

\begin{document}

\kaishu
\begin{titlepage}
        \heiti
        \vspace*{64pt}
        \begin{center}
            \fontsize{72pt}{0} 武汉纺织大学\\
            \vspace*{36pt}
            \fontsize{48pt}{0}{课\quad 程\quad 大\quad 作\quad 业\quad 模\quad 板}\\
            
            \vspace*{48pt}
            \LARGE（20\textbf{XX}\~{}20\textbf{XX}学年度\quad 第\textbf{X}学期）\\
            \vspace*{48pt}
            
            \LARGE { 面向多项式乘法问题的快速傅里叶变换算法设计与分析}\\
        	
        	\vspace*{100pt}
        	
            \Large 课程名称\ \ \underline{\makebox[200pt]{XXXXXXXXXXX}}\\
            \Large 作业名称\ \ \underline{\makebox[200pt]{XXXXXX}}\\
            \Large 学\qquad 院\ \ \underline{\makebox[200pt]{计算机与人工智能学院}}\\
            \Large 姓\qquad 名\ \ \underline{\makebox[200pt]{XXXXX}}\\
            \Large 班\qquad 级\ \ \underline{\makebox[200pt]{XXXXXXX班}}\\
            \Large 学\qquad 号\ \ \underline{\makebox[200pt]{XXXXXXXXX}}\\
            \Large 提交时间\ \ \underline{\makebox[200pt]{20\textbf{XX}年\textbf{XX}月\textbf{XX}日}}\\

            
            
%            \Large 地点\ \ \underline{\makebox[150pt]{}}\ \ 教师\ \ \underline{\makebox[108pt]{}}\\
        \end{center}
    \end{titlepage}

\newpage

\title{面向多项式乘法问题的快速傅里叶变换算法设计与分析\\
{\footnotesize}
\thanks{}
}

\author{\IEEEauthorblockN{XXXX\textsuperscript{1*}
}\\
\IEEEauthorblockA{\textit{\textsuperscript{1} 计算机与人工智能学院, 武汉纺织大学, 武汉, 中国 } \\
\textsuperscript{*}学生邮箱:  (xxxxxxx@wtu.edu.cn)}
}


\maketitle
\thispagestyle{fancy}
\chead{}
\renewcommand{\headrulewidth}{0pt}
\renewcommand{\footrulewidth}{0pt}
\pagestyle{fancy}
%\rfoot{\thepage}

\begin{abstract}
\\摘要：\\
\indent 散列表是一种数据结构，\\
\end{abstract}

\begin{IEEEkeywords}关键字：
IP 地址, 后继查找, van Emde Boas 树, 数据结构
\end{IEEEkeywords}


\section{引言}

文献\cite{AMaratea}中的字典问题出现在很多现实场景中，




\section{相关工作}



\subsection{字典问题}

\subsection{字典问题}

\subsection{字典问题}

如图\ref{figureDBOMotivationRel}所示，
\begin{figure}[ht]
\centerline{\includegraphics[width=3.5in]{figure_eps/散列表.png}}
\caption{散列链表.}
\label{figureDBOMotivationRel}
\end{figure}


\subsection{xxx}
\subsection{xxx}
\subsection{xxxxx}

\section{完全散列算法}

\subsection{动机}


\subsection{完全散列算法}
\subsubsection{全域散列簇}


\subsubsection{二级散列表的数据结构}

\subsubsection{完全散列算法}

\indent 如算法\ref{algProto-vEB-SUCCESSOR}所示，

\renewcommand{\algorithmicrequire}{\textbf{Input:}}
\renewcommand{\algorithmicensure}{\textbf{Output:}}
\begin{algorithm}[H]
\scriptsize   
\caption{SUCCESSOR(V,x)}           
\label{algProto-vEB-SUCCESSOR}                        
\begin{algorithmic}%[1]
\STATE /* 查找后继 */
\REQUIRE $V$表示vEB树,$x$表示插入的元素.
\STATE PROTO-vEB-SUCCESSOR(V,x)
\IF{$V.u==2$}
   \IF{$x==0$ and $V.A[1]==1$}
   \STATE return 1
   \ELSE
   \STATE return NIL
   \ENDIF
\ELSE
   \STATE offset=Proto-vEB-SUCCESSOR($V.cluster[high(x)]$,$low(x)$)
   \IF{$offset\neq NIL$}
      \STATE return index($high(x)$,offset)
   \ELSE
       \STATE succ-cluster=Proto-vEB-SUCCESSOR(V.summary,$high(x)$)
      \IF{$succ-cluster==-1$}
         \STATE return NIL
      \ELSE
         \STATE offset=Proto-vEB-MINIMUM(V.cluster[succ-cluster])
         \STATE return index(succ-cluster,offset)
      \ENDIF
   \ENDIF   
\ENDIF
\end{algorithmic}
\end{algorithm}



\subsection{理论分析}

\indent 如公式 Eq.~\ref{eqMIN}所示， \\

\begin{equation} T(u)=T(\sqrt{u})+O(1)
\label{eqMIN} 
\end{equation}





\section{实验结果}



\subsection{实验设计以及源代码对应函数}



\subsection{对比}


\subsection{讨论} 



\section{总结}





%\section*{Acknowledgment}


\bibliographystyle{IEEEtran}
% argument is your BibTeX string definitions and bibliography database(s)
\bibliography{IEEEabrv,refbib}


%\begin{thebibliography}{00}
%\bibitem{b1} G. Eason, B. Noble, and I. N. Sneddon, ``On certain integrals of Lipschitz-Hankel type involving products of Bessel functions,'' Phil. Trans. Roy. Soc. London, vol. A247, pp. 529--551, April 1955.
%\bibitem{b2} J. Clerk Maxwell, A Treatise on Electricity and Magnetism, 3rd ed., vol. 2. Oxford: Clarendon, 1892, pp.68--73.
%\bibitem{b3} I. S. Jacobs and C. P. Bean, ``Fine particles, thin films and exchange anisotropy,'' in Magnetism, vol. III, G. T. Rado and H. Suhl, Eds. New York: Academic, 1963, pp. 271--350.
%\bibitem{b4} K. Elissa, ``Title of paper if known,'' unpublished.
%\bibitem{b5} R. Nicole, ``Title of paper with only first word capitalized,'' J. Name Stand. Abbrev., in press.
%\bibitem{b6} Y. Yorozu, M. Hirano, K. Oka, and Y. Tagawa, ``Electron spectroscopy studies on magneto-optical media and plastic substrate interface,'' IEEE Transl. J. Magn. Japan, vol. 2, pp. 740--741, August 1987 [Digests 9th Annual Conf. Magnetics Japan, p. 301, 1982].
%\bibitem{b7} M. Young, The Technical Writer's Handbook. Mill Valley, CA: University Science, 1989.
%\end{thebibliography}

\end{document}
